\begin{resumo}

JUNIOR, Edson Ramos da Silva; MOTTA, Bernardo Cordeiro. \textbf{Plataforma Web Escalável e Segura para Venda de Ingressos Online: Arquitetura, Requisitos e Implementação}. 2025. Plano de Projeto — Curso de Análise e Desenvolvimento de Sistemas, Centro Universitário de Viçosa, Viçosa, 2025.

A industria global de eventos enfrenta desafios significativos na comercializacao digital de ingressos, incluindo problemas de escalabilidade durante picos de demanda, vulnerabilidades de seguranca e experiencias de usuario inadequadas. Este trabalho apresenta o projeto arquitetural e a implementacao de uma plataforma web robusta para venda de ingressos online, fundamentada em principios de engenharia de software e centrada nas necessidades dos stakeholders. A metodologia adotada incluiu analise detalhada dos stakeholders atraves de personas e mapa de empatia, seguida pela especificacao completa de requisitos funcionais e nao funcionais. A solucao implementada contempla uma arquitetura full-stack baseada em Next.js e MongoDB, sistema de autenticacao seguro com papeis por evento, gestao completa de eventos com tipos de ingresso e lotes, busca unificada com filtros, fluxo de compra com pagamento simulado e taxa de servico de 5\%, geracao de ingressos digitais nominais com QR Codes dinamicos assinados por HMAC, controle de acesso (check-in) com funcionamento offline e painel de analytics com intervalos temporais e comparativos de crescimento. Em conformidade com o carater academico do trabalho, o modulo de pagamento e simulado, sem integracao a um gateway financeiro real. O projeto adota como referencia padroes de seguranca como o OWASP Top 10, aderencia a LGPD e diretrizes de acessibilidade WCAG 2.1, e a analise de riscos e restricoes garante a viabilidade tecnica e operacional da solucao.

\textbf{Palavras-chave:} sistemas de ingressos; arquitetura full-stack; Next.js; MongoDB; seguranca da informacao; escalabilidade.

\end{resumo}
